\documentclass[11pt,a4paper]{scrartcl}

\usepackage[utf8]{inputenc}
\usepackage[T1]{fontenc}
\usepackage[margin=2.4cm]{geometry}
\usepackage{amsmath,amssymb}
\usepackage{siunitx}
\usepackage{booktabs}
\usepackage{longtable}
\usepackage{enumitem}
\usepackage{xcolor}
\usepackage{tcolorbox}
\usepackage[hidelinks]{hyperref}

\setlist{itemsep=2pt, topsep=4pt}

% deutsche Anfuehrungszeichen auch ohne babel-german verfuegbar
\providecommand{\glqq}{\quotedblbase}
\providecommand{\grqq}{\textquotedblleft}

% deutsche Bezeichner auch ohne babel-german
\renewcommand{\contentsname}{Inhaltsverzeichnis}
\renewcommand{\tablename}{Tabelle}
\renewcommand{\figurename}{Abbildung}

\definecolor{kastenrand}{RGB}{60,90,140}

\newtcolorbox{merke}[1][]{
  colback=blue!3, colframe=kastenrand, boxrule=0.6pt,
  left=6pt, right=6pt, top=5pt, bottom=5pt, #1}

\newcommand{\Tm}{T_\mathrm{m}}
\newcommand{\Tamb}{T_\mathrm{amb}}
\newcommand{\Cm}{C_\mathrm{m}}

\title{Projektvorgehen\\[2mm]\large Modellierung und Simulation eines Photovoltaikmoduls}
\subtitle{VU Kontinuierliche Simulation, 325.118 -- Arbeitsgrundlage f\"ur die Gruppe}
\date{}
\author{}

\begin{document}

\maketitle
\vspace{-12mm}
\tableofcontents
\bigskip

\begin{merke}
\textbf{Zweck dieses Dokuments.} Es legt fest, wer was macht, in welcher
Reihenfolge gearbeitet wird und nach welchen Regeln Code und Protokoll
entstehen. Wer eine Regel brechen will, darf das, muss es aber ansagen und
im Entscheidungslog vermerken. Das Ziel ist, dass am Ende jeder von uns
jeden Teil erkl\"aren kann, auch den, den er nicht selbst geschrieben hat.
\end{merke}

\section{Was verlangt wird}

Die Angabe nennt acht Punkte:

\begin{enumerate}[label=\arabic*.]
  \item Modellbildung
  \item Implementierung der Systemgleichungen in MATLAB
  \item Literaturrecherche der Modellparameter
  \item Validierung des Modells
  \item Anwendungsfall
  \item Sensitivit\"atsanalyse
  \item Implementierung in Simulink
  \item Protokoll und Pr\"asentation
\end{enumerate}

Entscheidend ist der Satz am Ende der Angabe: jeder Schritt muss
begr\"undbar und nachvollziehbar sein. Bewertet wird also nicht in erster
Linie, ob die Kurve stimmt, sondern ob wir sagen k\"onnen, warum wir sie so
gerechnet haben. Das Protokoll ist entsprechend gewichtet: das \glqq Warum\grqq{}
steht \"uber dem \glqq Was\grqq.

\section{Das Modell in Kurzform}

Bilanzraum ist das Modul ohne Rahmen, betrachtet als ein einziger
Temperaturknoten. Aus dem ersten Hauptsatz folgt

\begin{equation}
  \Cm \frac{\mathrm{d}\Tm}{\mathrm{d}t}
    = Q_\mathrm{solar} - W_\mathrm{el} - Q_\mathrm{konv} - Q_\mathrm{rad}.
  \label{eq:bilanz}
\end{equation}

Das Modul hat genau einen unabh\"angigen Energiespeicher, seine
W\"armekapazit\"at. Der Zustandsvektor hat damit die Ordnung eins, und zwar
nicht weil wir das so gew\"ahlt haben, sondern weil die Physik es vorgibt.
Dieser Satz geh\"ort in die Pr\"asentation, er ist der sauberste Bezug zur
Vorlesung.

Die elektrische Leistung h\"angt selbst von $\Tm$ ab,

\begin{equation}
  W_\mathrm{el} = G\,A\,(\tau\alpha)\,\eta_\mathrm{ref}
     \left[1 - \beta_\mathrm{ref}(\Tm - T_\mathrm{ref})\right].
\end{equation}

Genau diese R\"uckkopplung macht aus der Bilanz eine Differentialgleichung.
Ohne sie w\"are $\Tm$ eine explizite Formel und das ganze Projekt w\"are
ein Taschenrechnerproblem.

\subsection*{Was wir bewusst weglassen}

Wir rechnen das \textbf{Basismodell}. Konvektion \"uber einen linearen Ansatz
$h = h_a + h_b v$, nicht \"uber die vier Nusselt-Str\"omungsregime von Tuncel
et al.\ mit Windrichtung. Gr\"unde:

\begin{itemize}
  \item Die Windrichtung aus den GeoSphere-Daten l\"asst sich ohne Kenntnis der
        genauen Modulausrichtung nicht sinnvoll nutzen.
  \item Wir k\"onnen die vier Regime mit unseren Daten nicht validieren, wir
        w\"urden also Komplexit\"at hinzuf\"ugen, die wir nicht pr\"ufen k\"onnen.
  \item Die Angabe verlangt ausdr\"ucklich \glqq sinnvolle Annahmen und
        Vereinfachungen\grqq. Weglassen mit Begr\"undung ist die Aufgabe, nicht
        deren Umgehung.
\end{itemize}

\begin{merke}
Wer den Nusselt-Ansatz sp\"ater doch nachr\"usten will: es ist genau eine
Datei zu \"andern, \texttt{calc\_h\_conv.m}. Der Rest des Codes bleibt
unber\"uhrt. So ist die Struktur gebaut.
\end{merke}

\section{Rollen}

\begin{longtable}{l p{0.62\textwidth}}
  \toprule
  Rolle & Zust\"andig f\"ur \\
  \midrule
  \endhead
  \textbf{A -- Parameter} &
    Literaturrecherche aller Modellparameter, F\"ullen von
    \texttt{init\_parameters.m}, Parametertabelle im Protokoll,
    Herleitung von $\Cm$ aus dem Schichtaufbau. \\
  \addlinespace
  \textbf{B -- Modellkern} &
    MATLAB-Kern (ODE-Funktion, Solverwahl, Wetterinterpolation),
    Validierung gegen die Paperdaten, Kapitel Numerik und Validierung. \\
  \addlinespace
  \textbf{C -- Daten und Anwendung} &
    GeoSphere-Daten beschaffen und aufbereiten, Anwendungsfall,
    Sensitivit\"atsanalyse, Kapitel 6 und 7. \\
  \addlinespace
  \textbf{D -- Simulink} &
    Simulink-Modell, Kopplung mit dem Batteriemodell, Vergleich gegen die
    MATLAB-L\"osung, Kapitel 8. \\
  \addlinespace
  \textbf{Alle} &
    Modellbildung in der Startsitzung, Annahmentabelle, Pr\"asentation,
    Gegenlesen der Kapitel der anderen. \\
  \bottomrule
\end{longtable}

\begin{merke}
\textbf{Der Modellkern wird von genau einer Person geschrieben.} Vier Leute
an derselben ODE-Funktion ist der schnellste Weg zu Code, den hinterher
niemand mehr versteht. Wer etwas am Kern braucht, sagt es Person B.
\end{merke}

Eine Person hat zus\"atzlich das Schlusswort \"uber die Struktur des
Protokolls. Das ist keine inhaltliche Autorit\"at, sondern verhindert nur,
dass am Ende zehn Kapitel in zehn verschiedenen Stilen dastehen.

\section{Reihenfolge und Meilensteine}

\begin{longtable}{l p{0.50\textwidth} l}
  \toprule
  Phase & Inhalt & Wer \\
  \midrule
  \endhead
  0 & Papers lesen, Bilanzgleichung gemeinsam auf Papier, Systemgrenze
      festlegen, erste Annahmenliste & alle \\
  0 & GeoSphere-Daten herunterladen (24.06.--01.07.2019, 10-Minuten-Werte) & C \\
  \addlinespace
  1 & Parameterrecherche, \texttt{init\_parameters.m} f\"ullen & A \\
  1 & MATLAB-Kern lauff\"ahig, Testlauf mit Platzhalterdaten & B \\
  \addlinespace
  2 & Paperdaten digitalisieren, Validierung, Fehlerma{\ss}e & B, A \\
  2 & GeoSphere-Daten einlesen, Anwendungsfall rechnen & C \\
  \addlinespace
  3 & Sensitivit\"atsanalyse & C \\
  3 & Simulink-Modell und Batterie & D \\
  \addlinespace
  4 & Protokoll ausformulieren, Abbildungen final, Pr\"asentation & alle \\
  \bottomrule
\end{longtable}

Phase 0 ist keine Formalie. Wenn die Systemgrenze und die Annahmen nicht
gemeinsam festgelegt sind, schreibt jeder sp\"ater ein anderes Modell auf.

\begin{merke}
\textbf{Datenbeschaffung sofort.} Alles ab Aufgabenpunkt~5 h\"angt an den
GeoSphere-Daten. Erfahrungsgem\"a{\ss} h\"angt es dort: Stationsauswahl,
Datenl\"ucken, Zeitzone (UTC oder MEZ), und ob Globalstrahlung horizontal
oder in Modulebene vorliegt. Diese vier Punkte kl\"aren, bevor jemand
anf\"angt zu rechnen.
\end{merke}

\section{Ordnerstruktur}

\begin{verbatim}
pv_projekt/
  matlab/
    init_parameters.m        alle Parameter, einzige Stelle mit Zahlen
    pv_thermal_ode.m         rechte Seite der Bilanz
    calc_w_el.m              elektrische Leistung
    calc_h_conv.m            Waermeuebergangskoeffizient
    calc_q_rad.m             Strahlungsverlust
    calc_errors.m            MAE, RMSE, MBE
    build_weather_struct.m   baut die Interpolanten
    load_weather_paper.m     Validierungsdaten
    load_weather_geosphere.m Messdaten Anwendungsfall
    run_validation.m         rechnet  -> results/validation.mat
    run_usecase.m            rechnet  -> results/usecase.mat
    run_sensitivity.m        rechnet  -> results/sensitivity.mat
    plot_validation.m        plottet  <- results/validation.mat
    plot_usecase.m
    plot_sensitivity.m
    fig_style.m              einheitliches Aussehen
    save_figure.m            Export als Vektor-PDF
    run_all.m                Steuerung ueber Flags
    data/  results/  figures/  simulink/
  latex/
    main.tex
    kapitel/01_einleitung.tex ... A_anhang.tex
    literatur.bib
  doku/
    Projektvorgehen.pdf      dieses Dokument
    annahmen.md              Entscheidungslog
\end{verbatim}

\section{Arbeitsregeln f\"ur den Code}

\begin{enumerate}[label=\textbf{R\arabic*.}, leftmargin=*]
  \item \textbf{Keine Zahlenwerte au{\ss}erhalb von \texttt{init\_parameters.m}.}
        Wer eine Zahl direkt in eine Rechnung schreibt, muss sie sp\"ater in
        drei Dateien suchen.
  \item \textbf{Rechnen und Plotten sind getrennt.} \texttt{run\_*.m}
        speichert nach \texttt{results/}, \texttt{plot\_*.m} l\"adt von dort.
        Wer nur die Achsenbeschriftung \"andert, startet keine Simulation neu.
  \item \textbf{Alles in SI, Temperaturen intern in Kelvin.} Umrechnung nach
        Grad Celsius erst beim Plotten. Einheitenfehler sind die h\"aufigste
        Fehlerquelle und die am schwersten zu findende.
  \item \textbf{Gleichungsnummern synchron halten.} Jede Gleichung im
        Protokoll bekommt ein \texttt{\textbackslash label}, im Code steht
        dar\"uber \texttt{\% Gl. (2.3)}. Kostet f\"unf Minuten und beantwortet
        in der Pr\"ufung jede \glqq wo steht das im Code\grqq-Frage sofort.
  \item \textbf{Sprechende Namen.} \texttt{T\_module} statt \texttt{x1},
        \texttt{h\_conv} statt \texttt{hc}.
  \item \textbf{Eine Datei, ein Zweck.} \"Uber etwa 100 Zeilen steckt meist
        eine Funktion drin, die ausgelagert geh\"ort.
  \item \textbf{Keine anonymen Funktionen mit Logik.} Ein Wrapper der Form
        \texttt{@(t,T) pv\_thermal\_ode(t,T,p,w)} ist in Ordnung, alles
        dar\"uber nicht.
  \item \textbf{Vektorisierung nur, wo sie sp\"urbar etwas bringt.} Eine
        lesbare Schleife schl\"agt einen cleveren Einzeiler. Wir werden auf
        Nachvollziehbarkeit bewertet, nicht auf Rechenzeit.
  \item \textbf{Kein Feature ohne Aufgabenbezug.} Wenn niemand sagen kann,
        welchen der acht Punkte eine Funktion bedient, fliegt sie raus.
  \item \textbf{Versionierung.} Git, oder zumindest datierte Ordner. Nicht
        \texttt{final\_final\_v3.m}.
\end{enumerate}

\subsection{Zwei technische Stolpersteine}

\textbf{Wetterdaten m\"ussen interpoliert werden.} Der ODE-Solver w\"ahlt
seine Zeitschritte selbst und fragt Zeitpunkte ab, die zwischen den
Messwerten liegen. Die Interpolanten werden einmal in
\texttt{build\_weather\_struct.m} erzeugt und in die Struct \texttt{w}
gelegt, nicht bei jedem Aufruf der rechten Seite neu. Lineare Interpolation,
weil ein Spline bei Bew\"olkungsspr\"ungen \"uberschwingen und negative
Einstrahlung erzeugen kann.

\textbf{Die Solverwahl muss begr\"undet werden.} Unser System ist skalar,
ein Steifheitsverh\"altnis ist damit gar nicht definiert, ein steifes
Problem kann es also nicht sein. Daraus folgt ein explizites Verfahren,
konkret \texttt{ode45} (Dormand-Prince, Ordnung 5 mit eingebetteter Ordnung 4
zur Schrittweitensteuerung). Toleranzen setzen wir bewusst auf
\num{1e-6} statt die Defaults zu nehmen, und wir zeigen einmal, dass sich
das Ergebnis bei strengerer Toleranz nicht mehr \"andert. Genau das ist der
Punkt, an dem die Vorlesung gepr\"uft wird, \glqq haben wir immer so gemacht\grqq{}
reicht nicht.

\section{Arbeitsregeln f\"ur das Protokoll}

Kapitelstruktur:

\begin{enumerate}[leftmargin=*]
  \item Einleitung und Aufgabenstellung
  \item Modellbildung (Systemabgrenzung, Bilanz, Zustandswahl, Teilmodelle,
        \textbf{Annahmentabelle})
  \item Modellparameter (Recherche, Parametertabelle mit Quellen)
  \item Numerische Umsetzung (Programmstruktur, Solverwahl, Interpolation)
  \item Validierung (Vorgehen, Fehlerma{\ss}e, Ergebnisse, Diskussion der Abweichungen)
  \item Anwendungsfall (Datengrundlage, Ergebnisse, dominanter Verlustterm)
  \item Sensitivit\"atsanalyse
  \item Simulink und Batterieladung
  \item Diskussion und Grenzen des Modells
  \item Verwendung von KI-Werkzeugen
\end{enumerate}

Jedes Kapitel ist in Overleaf eine eigene Datei unter \texttt{kapitel/}.
So k\"onnen vier Leute gleichzeitig schreiben, ohne sich zu blockieren. Ganz
oben in jeder Datei steht ein Kommentarblock mit \texttt{VERANTWORTLICH:}
und einer Liste, was in dieses Kapitel geh\"ort und was nicht.

\subsection{Die Annahmentabelle}

Das ist das St\"uck, das \"uber die Note entscheidet. Drei Spalten: was wurde
angenommen, warum ist das zul\"assig, was kostet es voraussichtlich. Beispiel:

\begin{center}
\small
\begin{tabular}{p{0.26\textwidth} p{0.34\textwidth} p{0.28\textwidth}}
  \toprule
  Annahme & Begr\"undung & Erwartete Auswirkung \\
  \midrule
  Modul als ein Temperaturknoten &
    Laminatdicke rund \SI{3}{\milli\metre}, Biot-Zahl klein &
    kein Gradient abbildbar, im Mittel vernachl\"assigbar \\
  Bodentemperatur gleich $\Tamb$ &
    keine Messdaten verf\"ugbar &
    untersch\"atzt R\"uckseitenverluste bei aufgeheiztem Untergrund \\
  \bottomrule
\end{tabular}
\end{center}

\begin{merke}
\textbf{Regel: neue Annahme im Code heisst neue Zeile in der Tabelle.} Wer
im Code etwas vereinfacht, ohne es einzutragen, hat die Arbeit nur
verschoben, nicht erledigt.
\end{merke}

Parallel dazu f\"uhren wir \texttt{doku/annahmen.md} als Entscheidungslog:
Datum, Entscheidung, Begr\"undung, wer. Eine Zeile pro Eintrag, keine Prosa.
Daraus wird sp\"ater die Tabelle.

\subsection{Abbildungen}

Die Protokolltipps der LVA zeigen konkrete Negativbeispiele. Was dort
kritisiert wird:

\begin{itemize}
  \item fehlende Achsenbeschriftung, unbekannte Einheiten
  \item Legende vorhanden, aber unsauber beschriftet (etwa \texttt{phi0})
  \item gleiche Linienart f\"ur alle Kurven, Unterscheidung nur \"uber Farbe
  \item zu kleine Plots, Schriftgr\"o{\ss}e passt nicht zum Flie{\ss}text
  \item Redundanz zwischen Plottitel, Achsenbeschriftung und Bildunterschrift
  \item Simulink-Scopes direkt ins Protokoll kopiert
  \item nichtssagende Bildunterschriften wie \glqq Simulationsgraph Simulation 2\grqq
\end{itemize}

Die zugeh\"orige Checkliste f\"ur wissenschaftliche Arbeiten verlangt
zus\"atzlich: Vektorgrafik, sinnvoller Bildausschnitt ohne gro{\ss}e leere
Fl\"achen, kein Informationsverlust im Schwarzwei{\ss}druck, fortlaufende
Nummerierung, und eine Bildunterschrift, die die Abbildung unabh\"angig vom
Text verst\"andlich macht und die Kernaussage benennt.

Praktisch hei{\ss}t das f\"ur uns: \texttt{fig\_style.m} am Anfang jedes
Plotskripts aufrufen, Export \"uber \texttt{save\_figure.m} als
Vektor-PDF, unterschiedliche Linienarten statt nur Farben, und meist gar
keinen Plottitel setzen, weil die Bildunterschrift das \"ubernimmt.

\subsection{KI-Werkzeuge}

Die Nutzung ist erlaubt und muss nicht verschleiert werden. Kapitel 10
beschreibt: welche Werkzeuge, wof\"ur konkret, und vor allem \emph{wie
gepr\"uft wurde}. Der letzte Punkt ist der, auf den es ankommt. Jede
\"ubernommene Formel wird gegen die Prim\"arquelle gepr\"uft, jeder Parameter
gegen die zitierte Arbeit. Was verworfen wurde, wird ebenfalls erw\"ahnt.

\section{Offene Punkte}

\begin{longtable}{p{0.52\textwidth} l l}
  \toprule
  Was & Wer & Bis wann \\
  \midrule
  \endhead
  GeoSphere-Datensatz ausw\"ahlen und herunterladen & C & \\
  Zeitzone und Datenl\"ucken kl\"aren & C & \\
  Paperdaten aus Tuncel Abb.~1 digitalisieren & B & \\
  Umgebungstemperatur f\"ur die Validierung festlegen (Paper macht keine Angabe) & B & \\
  Alle \texttt{TODO Quelle} in \texttt{init\_parameters.m} belegen & A & \\
  Schichtaufbau f\"ur $\Cm$ mit Quellen hinterlegen & A & \\
  Modulausrichtung und Neigung festlegen & alle & \\
  Batteriemodell von der LVA anfordern & D & \\
  Simulink-Onramp durcharbeiten & D & \\
  \bottomrule
\end{longtable}

\section{Woran wir merken, dass es gut l\"auft}

\begin{itemize}
  \item Jede Zahl im Code hat eine Quelle im Kommentar.
  \item Jede Gleichung im Protokoll hat eine Nummer, die im Code vorkommt.
  \item Die Annahmentabelle w\"achst mit dem Code mit, nicht erst am Schluss.
  \item Wer ein Kapitel nicht geschrieben hat, kann es trotzdem erkl\"aren.
  \item Ein \texttt{plot\_*.m}-Aufruf braucht keine Neuberechnung.
\end{itemize}

\end{document}
